% GENERATED FILE. Edit canonical lessons, then run npm run generate:books.
% canonical-source-hash: fnv1a64:64054183b1360378
% canonical-lessons: GE-C292-warnen, GE-C292-rettung, GE-C292-bergen, GE-C292-evakuieren, GE-C292-ueberleben

\chapter{Die Rettung}
\label{ch:ge-die-rettung}
\hlchaptermodality{292}

\begin{chapteropening}
\textbf{By the end of this chapter:} \emph{I can report a rescue: to warn, a rescue, to bring out alive, to evacuate and to survive.}

Complete the last lesson of chapter 292: I can report a rescue: to warn, a rescue, to bring out alive, to evacuate and to survive.
\end{chapteropening}

\section[warnen]{warnen — to warn}

\label{lesson:GE-C292-warnen}



\begin{quote}
Before the new one: say the German for the paramedic, then the German for the emergency room.
\end{quote}

\subsection*{You'll want to know: warnen}
\textbf{warnen} — \textquotedblleft{}to warn\textquotedblright{}.

\begin{quote}
Die Nachbarin warnte uns vor dem Hund.
\end{quote}

\begin{quote}
\emph{The neighbour warned us about the dog.}
\end{quote}

\emph{Vor etwas warnen} — to warn about something. \emph{Die Warnung} — the warning.

\begin{cousinweb}
\emph{warnen} is inherited from Middle High German \emph{warnen} \textquotedblleft{}to watch over, to protect\textquotedblright{} (EtymDB), a cousin of English \emph{warn}.
\end{cousinweb}

\begin{grammarlens}[title={What you've built}]
One more word for telling, explaining and describing at B1.
\end{grammarlens}

\subsection*{Your turn}
\begin{itemize}
\raggedright
  \item \emph{Say it:} \emph{warnen}
  \item \emph{Say it:} \emph{warnen}, once more
  \item \emph{Say it:} a sentence of your own with \emph{warnen}
  \item \emph{Recall:} say the German for the paramedic, then the German for the emergency room, then say \emph{warnen} again
\end{itemize}

\begin{tcolorbox}[breakable,colback=teal!4,colframe=teal!35!black,title={Before you move on}]
What does warnen mean? (\textquotedblleft{}To warn\textquotedblright{}.) Say it once more.
\end{tcolorbox}
\section[die Rettung]{die Rettung — the rescue}

\label{lesson:GE-C292-rettung}



\begin{quote}
Before the new one: say the German for the emergency room, then the German for to warn.
\end{quote}

\subsection*{You'll want to know: die Rettung}
\textbf{die Rettung} — \textquotedblleft{}the rescue\textquotedblright{}.

\begin{quote}
Die Rettung kam in letzter Minute.
\end{quote}

\begin{quote}
\emph{The rescue came at the last minute.}
\end{quote}

Plural: \emph{die Rettungen}. In Austria, \emph{die Rettung} is also the ambulance service.

\begin{cousinweb}
\emph{Rettung} comes from Middle High German \emph{rettunge} (EtymDB), built on \emph{retten} \textquotedblleft{}to save\textquotedblright{}, with the ending \emph{-ung}.
\end{cousinweb}

\begin{grammarlens}[title={What you've built}]
One more word for telling, explaining and describing at B1.
\end{grammarlens}

\subsection*{Your turn}
\begin{itemize}
\raggedright
  \item \emph{Say it:} \emph{die Rettung}
  \item \emph{Say it:} \emph{die Rettung}, once more
  \item \emph{Say it:} a sentence of your own with \emph{die Rettung}
  \item \emph{Recall:} say the German for the emergency room, then the German for to warn, then say \emph{die Rettung} again
\end{itemize}

\begin{tcolorbox}[breakable,colback=teal!4,colframe=teal!35!black,title={Before you move on}]
What does die Rettung mean? (\textquotedblleft{}The rescue\textquotedblright{}.) Say it once more.
\end{tcolorbox}
\section[bergen]{bergen — to bring out alive, to recover}

\label{lesson:GE-C292-bergen}



\begin{quote}
Before the new one: say the German for to warn, then the German for the rescue.
\end{quote}

\subsection*{You'll want to know: bergen}
\textbf{bergen} — \textquotedblleft{}to bring out alive, to recover\textquotedblright{}.

\begin{quote}
Die Feuerwehr barg zwei Menschen aus dem Auto.
\end{quote}

\begin{quote}
\emph{The fire brigade brought two people out of the car.}
\end{quote}

Strong verb: \emph{barg}, \emph{hat geborgen}. \emph{Geborgen} also means 'safe and sheltered'.

\begin{cousinweb}
\emph{bergen} is inherited from Old High German \emph{bergan}, which EtymDB traces to Proto-Germanic \emph{berganą} and a root meaning \textquotedblleft{}to protect\textquotedblright{}.
\end{cousinweb}

\begin{grammarlens}[title={What you've built}]
One more word for telling, explaining and describing at B1.
\end{grammarlens}

\subsection*{Your turn}
\begin{itemize}
\raggedright
  \item \emph{Say it:} \emph{bergen}
  \item \emph{Say it:} \emph{bergen}, once more
  \item \emph{Say it:} a sentence of your own with \emph{bergen}
  \item \emph{Recall:} say the German for to warn, then the German for the rescue, then say \emph{bergen} again
\end{itemize}

\begin{tcolorbox}[breakable,colback=teal!4,colframe=teal!35!black,title={Before you move on}]
What does bergen mean? (\textquotedblleft{}To bring out alive, to recover\textquotedblright{}.) Say it once more.
\end{tcolorbox}
\section[evakuieren]{evakuieren — to evacuate}

\label{lesson:GE-C292-evakuieren}



\begin{quote}
Before the new one: say the German for the rescue, then the German for to bring out alive.
\end{quote}

\subsection*{You'll want to know: evakuieren}
\textbf{evakuieren} — \textquotedblleft{}to evacuate\textquotedblright{}.

\begin{quote}
Wegen eines Feuers wurde die Schule evakuiert.
\end{quote}

\begin{quote}
\emph{Because of a fire the school was evacuated.}
\end{quote}

Regular, without \emph{ge-}: \emph{hat evakuiert}. Stress: \emph{e-va-ku-IE-ren}.

\begin{cousinweb}
\emph{evakuieren} is a learned borrowing from Latin \emph{ēvacuāre} \textquotedblleft{}to empty out\textquotedblright{} (EtymDB), built on \emph{vacuus} \textquotedblleft{}empty\textquotedblright{} (Lewis \& Short).
\end{cousinweb}

\begin{grammarlens}[title={What you've built}]
One more word for telling, explaining and describing at B1.
\end{grammarlens}

\subsection*{Your turn}
\begin{itemize}
\raggedright
  \item \emph{Say it:} \emph{evakuieren}
  \item \emph{Say it:} \emph{evakuieren}, once more
  \item \emph{Say it:} a sentence of your own with \emph{evakuieren}
  \item \emph{Recall:} say the German for the rescue, then the German for to bring out alive, then say \emph{evakuieren} again
\end{itemize}

\begin{tcolorbox}[breakable,colback=teal!4,colframe=teal!35!black,title={Before you move on}]
What does evakuieren mean? (\textquotedblleft{}To evacuate\textquotedblright{}.) Say it once more.
\end{tcolorbox}
\section[überleben]{überleben — to survive}

\label{lesson:GE-C292-ueberleben}



\begin{quote}
Before the new one: say the German for to bring out alive, then the German for to evacuate.
\end{quote}

\subsection*{You'll want to know: überleben}
\textbf{überleben} — \textquotedblleft{}to survive\textquotedblright{}.

\begin{quote}
Alle im Bus überlebten den Unfall.
\end{quote}

\begin{quote}
\emph{Everyone on the bus survived the accident.}
\end{quote}

Not separable: \emph{er überlebt}, \emph{er hat überlebt}.

\begin{cousinweb}
\emph{überleben} joins \emph{über} \textquotedblleft{}over, beyond\textquotedblright{} and \emph{leben} \textquotedblleft{}to live\textquotedblright{}, from Proto-Germanic \emph{libjaną} (EtymDB): to live beyond something.
\end{cousinweb}

\begin{grammarlens}[title={What you've built}]
One more word for telling, explaining and describing at B1.
\end{grammarlens}

\subsection*{Your turn}
\begin{itemize}
\raggedright
  \item \emph{Say it:} \emph{überleben}
  \item \emph{Say it:} \emph{überleben}, once more
  \item \emph{Say it:} a sentence of your own with \emph{überleben}
  \item \emph{Recall:} say the German for to bring out alive, then the German for to evacuate, then say \emph{überleben} again
\end{itemize}

\begin{tcolorbox}[breakable,colback=teal!4,colframe=teal!35!black,title={Before you move on}]
What does überleben mean? (\textquotedblleft{}To survive\textquotedblright{}.) Say it once more.
\end{tcolorbox}
